% Vista científica exacta materializada desde una rama condicional.
% Fuente: source/demostraciones_primarias/31_06c_sectores_topologicos_no_omision.tex; rama: HMTIncludeCMBlock.
% No modifica el contenido matemático de la rama de origen.
\chapter[Structure complexe \texorpdfstring{\(\widetilde J_M^{\,2}=-I\)}{J M au carré = -I} et secteurs topologiques]
{Structure complexe \texorpdfstring{\(\widetilde J_M^{\,2}=-I\)}{J M au carré = -I} sur la strate fixe, involutions dodécaphasiques et secteurs topologiques}
\label{chap:sectores-topologicos-no-omision}
\index{secteur topologique}\index{Fibonacci, anneau basé}
\index{Majorana}\index{Ising}\index{groupe de tresses}

L’involution \(C_M\), l’anneau basé associé à \(F_{\rm av}\), les caractères de tresse et les catégories pointées appartiennent à des strates algébriques distinctes. Ils sont présentés avec leurs domaines et leurs contrôles d’identification afin qu’une coïncidence de nom ne remplace pas le morphisme qui devrait les relier. Le dénombrement de l’atlas, la hiérarchie spectrale infinie dont les huit hauteurs historiques sont les témoins initiaux, et les complétions statistiques conservent les démonstrations de \cref{thm:censo-atlas-81,thm:descenso-necesario-fav,%
thm:principio-exclusion-pauli,thm:dos-cuatro-pi-estadistica}.

Chaque résultat énonce directement son domaine, sa construction et sa conclusion ; aucune coïncidence de nom ne modifie ces dépendances.

\section{Strate périodique de l’involution dodécaphasique}
\label{sec:involucion-cm-censo}

Le domaine \(\mathscr T_{12}^{\rm word}\), la conjugaison \(C_M\mathbf t=-\operatorname{rev}(\mathbf t)\), son caractère involutif et le dénombrement \(\lvert\operatorname{Fix}(C_M)\rvert=3^6=729\) ont leur emplacement principal dans les \cref{def:palabras-cm,thm:involucion-cm-censo}, au sein de la bande de particule. Seule est ajoutée ici la condition périodique nécessaire à la strate topologique. Si \(r\) désigne le décalage cyclique \((r\mathbf t)_m=t_{m-1}\), nous posons
\[
 \operatorname{Per}_6
 :=\{\mathbf t\in\mathscr T_{12}^{\rm word}:r^6\mathbf t=\mathbf t\}.
\]

\begin{proposition}[Dénombrement de la strate périodique]
\label{prop:estrato-periodico-cm-rev7} En imposant une périodicité de période divisant six,
\[
 \bigl|\operatorname{Fix}(C_M)\cap\operatorname{Per}_6\bigr|
 =3^3=27.
\]
Cette identité est exacte.
\end{proposition}

\begin{proof}
Le \cref{thm:involucion-cm-censo} donne la forme générale des points fixes. Si, de plus, \(r^6\mathbf t=\mathbf t\), la seconde moitié du mot coïncide avec la première. En écrivant les six premières coordonnées, les deux conditions équivalent à
\[
 (t_0,t_1,t_2,t_3,t_4,t_5)
 =(a,b,c,-c,-b,-a).
\]
Les paramètres \(a,b,c\) sont indépendants et prennent trois valeurs ; l’intersection possède donc \(3^3\) éléments.
\end{proof}

La forme générale d’un point fixe, avant d’exiger la période six, est
\[
 (a,b,c,d,e,f,-f,-e,-d,-c,-b,-a).
\]
Cette forme en coordonnées permet de vérifier manuellement l’involution et le dénombrement, et évite d’interpréter l’entier \(729\) au moyen d’un autre objet homonyme.

\begin{remark}[Dénombrement exhaustif de l’involution]
\label{rem:procedencia-n40-cm} L’énumération des \(3^{12}=531\,441\) mots ne relève aucune violation de l’involution et donne
\[
 |\operatorname{Fix}(C_M)|=729,
 \qquad
 |\operatorname{Fix}(C_M)\cap\operatorname{Per}_6|=27.
\]
La démonstration principale et la \cref{prop:estrato-periodico-cm-rev7} expliquent ces deux nombres indépendamment de l’énumération : les points fixes ont six coordonnées libres, et la condition supplémentaire de période six réduit le choix à trois. Le dénombrement exhaustif confirme qu’il n’existe aucune orbite exceptionnelle en dehors de cette description.
\end{remark}

\begin{criterion}[Critères de réfutation du dénombrement \(C_M\)] \label{crit:falsadores-cm} Chacun des faits suivants réfuterait le résultat : un mot tel que \(C_M^2\mathbf t\ne\mathbf t\) ; un point fixe ne possédant pas la forme en coordonnées précédente ; un total exhaustif différent de \(729\) ; ou un total différent de \(27\) après avoir imposé \(r^6\mathbf t=\mathbf t\). La coïncidence de ce \(729\) avec un autre cardinal n’identifie pas leurs domaines.
\end{criterion}

\section{Structure quadratique transportée sur la strate fixe}
\label{sec:estructura-cuadratica-estrato-fijo-cm}

Nous identifions l’alphabet des signes au corps ternaire au moyen de
\[
 \{-1,0,+1\}\xrightarrow{\sim}\mathbb F_3,
 \qquad
 -1\longmapsto2,\quad 0\longmapsto0,\quad +1\longmapsto1.
\]
Avec cette identification, \(\mathscr T_{12}^{\rm word}\cong\mathbb F_3^{12}\), et tant \(C_M=-\operatorname{rev}\) que le décalage \(r\) sont des opérateurs \(\mathbb F_3\)-linéaires. Par conséquent, \(\operatorname{Fix}(C_M)\) et \(\operatorname{Per}_6=\ker(r^6-I)\) sont des sous-espaces vectoriels.

La forme en coordonnées des points fixes définit l’isomorphisme linéaire
\[
 s_{\rm fix}:\mathbb F_3^6\xrightarrow{\sim}\operatorname{Fix}(C_M),
\]
\[
 s_{\rm fix}(a,b,c,d,e,f)
 =(a,b,c,d,e,f,-f,-e,-d,-c,-b,-a).
\]
Soit \(J_W\) l’opérateur de Paley--Witt construit dans le \cref{thm:f9-tres-desde-paley}, avec la base, la polarisation et l’orientation qui y sont déclarées.

\begin{proposition}[Structure complexe \(\widetilde J_M^{\,2}=-I\) sur la strate fixe]
\label{prop:raiz-menos-uno-estrato-fijo-cm} Le transport
\[
\widetilde J_M:=s_{\rm fix}J_Ws_{\rm fix}^{-1}
\in\operatorname{End}_{\mathbb F_3}
   \bigl(\operatorname{Fix}(C_M)\bigr)
\]
satisfait
\[
 \widetilde J_M^{\,2}=-I.
\]
\end{proposition}

\begin{proof}
Par conjugaison,
\[
 \widetilde J_M^{\,2}
 =s_{\rm fix}J_Ws_{\rm fix}^{-1}
  s_{\rm fix}J_Ws_{\rm fix}^{-1}
 =s_{\rm fix}J_W^2s_{\rm fix}^{-1}
 =-I.
\]
\end{proof}

Le résultat est démontré avec la carte \(s_{\rm fix}\), la base, la polarisation et l’orientation comme structure initiale explicite. Il n’identifie pas \(C_M\) à \(J_W\) : le premier est une involution de l’espace des mots et le second une structure quadratique sur sa strate fixe. Il ne démontre pas non plus que le transport soit naturel relativement à toute la dynamique TPK.

Posons
\[
 W_{27}
 :=\operatorname{Fix}(C_M)\cap\operatorname{Per}_6.
\]
L’inclusion
\[
 W_{27}\hookrightarrow\operatorname{Fix}(C_M)
\]
est canonique, et la \cref{prop:estrato-periodico-cm-rev7} prouve que \(\dim_{\mathbb F_3}W_{27}=3\).

\begin{corollary}[Obstruction quadratique dans la strate à vingt-sept éléments]
\label{cor:obstruccion-cuadratica-w27} Le sous-espace \(W_{27}\) ne peut être stable sous \(\widetilde J_M\).
\end{corollary}

\begin{proof}
S’il était stable, la relation \(\widetilde J_M^2=-I\) en ferait un module sur
\[
 \mathbb F_3[X]/(X^2+1)\cong\mathbb F_9.
\]
Tout espace vectoriel sur \(\mathbb F_9\) est de dimension paire lorsqu’il est considéré sur \(\mathbb F_3\), ce qui contredit \(\dim_{\mathbb F_3}W_{27}=3\).
\end{proof}

L’égalité de cardinal entre \(W_{27}\) et les vingt-sept droites de Schläfli ne fournit pas d’application entre ces deux objets. Une telle relation exigerait la construction d’une application d’incidence préservant l’orientation centrale et l’action du stabilisateur, ainsi qu’une clôture dynamique compatible. L’atlas de \(81\) cellules ne s’interpose pas entre \(W_{27}\) et la strate fixe de \(729\) états : l’atlas des particules appartient à un autre domaine.

\section[Niveaux de Majorana]{Stratification mathématique de l’objet de Majorana en quatre niveaux}
\label{sec:cuatro-niveles-majorana}

L’inversion \(C_M\), la parité extérieure, un fermion relativiste de Majorana et un mode zéro de Majorana sont des objets de types distincts :

\begin{enumerate}
 \item \(C_M=-\operatorname{rev}\) agit sur \(\mathscr T_{12}^{\rm word}\). Son involutivité et son dénombrement sont démontrés exactement ; cette construction ne fournit à elle seule ni spineur, ni équation de champ, ni terme de masse.

 \item Le caractère de signe \(q=-1\) et la complétion extérieure appartiennent à l’algèbre d’échange. Relativement aux CAR et à la graduation déclarée, ils produisent \(N_s^2=N_s\), le spectre \(\{0,1\}\) et le retour d’orientation \(2\pi/4\pi\). Ces identités se déduisent des CAR et de la graduation ; elles ne se déduisent pas du seul dénombrement de \(C_M\).

 \item Un fermion relativiste de Majorana est un objet spinoriel muni d’une représentation relativiste, d’une conjugaison de charge compatible, d’une condition d’autoconjugaison, d’un terme de masse et d’une dynamique. Ces données ne font pas partie de la définition de \(C_M\) et ne se déduisent pas de son dénombrement.

 \item Un mode zéro de Majorana, ou, de manière équivalente, le secteur non abélien d’une théorie d’Ising au sens topologique, requiert des objets simples \(\mathbf 1,\psi,\sigma\) avec
 \[
  \psi\otimes\psi\simeq\mathbf 1,\qquad
  \psi\otimes\sigma\simeq\sigma,\qquad
  \sigma\otimes\sigma\simeq\mathbf 1\oplus\psi,\qquad
  d_\sigma=\sqrt2,
 \]
 ainsi qu’un associateur et un tressage cohérents. Le gap, le dédoublement et la plateforme relèvent d’une réalisation physique supplémentaire ; les règles de fusion précédentes ne définissent pas une masse de pôle universelle.
\end{enumerate}

La séparation concerne les domaines et non la nomenclature : involution de mots, parité CAR, spineur autoconjugué et objet simple non abélien sont quatre objets mathématiques distincts. En l’absence d’un morphisme explicite entre eux, aucun n’est identifié à un autre du seul fait qu’ils portent le nom de « Majorana ».
